\section{Separate success values from error positions}\label{expression-results}
The meaning of \texttt{v1} depends on \texttt{v0}. On success it is a signed
integer result; on failure it is an unsigned, nonnegative byte offset no greater
than the supplied length. Inspect the status before interpreting the second result.
The value minus two in \texttt{v0}, for example, is an overflow status rather than
an expression result. A successful expression can itself evaluate to minus two,
which is then returned in \texttt{v1} with a zero status.

A syntax error points at the unexpected byte, or at end of input if a required
operand or closing parenthesis is absent. A decimal-literal overflow points at
the first digit whose inclusion makes that literal exceed its maximum. Overflow
from a unary operation points at its sign byte. Overflow or division by zero from
a binary operation points at that operator's byte. These positions describe the
operation that failed; they need not be the rightmost byte that was examined.

\begin{table}[htbp]
\centering
\caption{Diagnostic examples. Positions are decimal byte offsets; all displayed inputs contain only ASCII bytes.}
\tagpdfsetup{table/header-rows={1}}
\begin{tabular}{lrrl}
\toprule
Input & Status & Offset & Reason\\
\midrule
\texttt{1 2} & $-1$ & 2 & Unexpected second literal\\
\texttt{(1} & $-1$ & 2 & Missing closing parenthesis\\
\texttt{1+} & $-1$ & 2 & Missing right operand\\
\texttt{2147483648} & $-2$ & 9 & Literal's final digit exceeds range\\
\texttt{2147483647+1} & $-2$ & 10 & Addition exceeds range\\
\texttt{1/(2-2)} & $-3$ & 1 & Division has a zero denominator\\
\bottomrule
\end{tabular}
\end{table}
